\subsubsection{Descenso de operaciones al cociente de evaluación}
\label{sec:ch-descenso-fibras}

El cociente anterior identifica historias que expresan el mismo valor, no
sus registros completos. El criterio siguiente caracteriza exactamente
cuándo una operación de historias induce una operación en ese cociente sin
escoger una historia privilegiada dentro de cada fibra; su formulación
continúa la construcción genealógica establecida en las secciones anteriores.

\begin{proposition}[Criterio de descenso por fibras de evaluación]
\label{prop:ch-descenso-fibras}
Sea \(q:\Omega\to J\) la evaluación sobre su imagen de una coordenatización
genealógica. Sea \(\star:D\to\Omega\) una operación con
\(D\subseteq\Omega^2\) saturado por las fibras de \(q\times q\).
Existe una única operación
\(\overline\star:(q\times q)(D)\to J\) tal que
\[
 q(x\star y)=q(x)\,\overline\star\,q(y)
\]
si y sólo si, para pares del dominio,
\[
 q(x)=q(x'),\quad q(y)=q(y')
 \quad\Longrightarrow\quad q(x\star y)=q(x'\star y').
\]
\end{proposition}

\begin{proof}
Si existe la operación visible, su resultado depende sólo de los dos
valores, lo que prueba la necesidad. Para la suficiencia, se define
\(s\,\overline\star\,t=q(x\star y)\), donde \(q(x)=s\) y
\(q(y)=t\). La saturación conserva el dominio al cambiar representantes;
la condición enunciada conserva el resultado. La sobreyectividad de
\(q\times q\) sobre la imagen del dominio prueba la unicidad.
\end{proof}

Una coordenada no reemplaza automáticamente una historia como argumento de
todas sus operaciones: el criterio determina cuándo esa sustitución es
legítima. En particular, transportar una operación a través de una biyección
no basta para identificarla con una operación nativa anterior a la lectura.
Si una coordenatización \(\Omega\) es compacta y \(q\) es continua y sobreyectiva
sobre un espacio Hausdorff \(J\), la evaluación induce además un
homeomorfismo \(\Omega/{\sim_q}\to J\): es una biyección continua de un
compacto a un Hausdorff. Es el \emph{cociente}, y no necesariamente el
espacio de historias, el que se identifica topológicamente con \(J\).
